\chapter[Decisión multicriterio y estudio de alternativas]{Análisis de decisión multicriterio y estudio cuantitativo de alternativas}\label{cap:decision}

Con los pesos base del equipo, el paracaídas y el autogiro quedan \emph{empatados} en el
AHP ($0{,}218$ frente a $0{,}217$), y TOPSIS prefiere el autogiro ($0{,}79$ frente a
$0{,}66$ del coaxial y $0{,}61$ del paracaídas). Este capítulo reemplaza la matriz de puntos de la
tabla~\ref{tab:pugh} por un estudio cuantitativo. Primero modela cada alternativa de la
figura~\ref{fig:alternativas} dentro de la envolvente del CanSat. Luego aplica dos
métodos de decisión (AHP y TOPSIS) y mide cuán robusta es la elección frente a los pesos,
los juicios y el método. El autogiro no gana por sus prestaciones, que son peores que las del paracaídas, sino
porque el proyecto valora su contenido técnico. En el AHP conviene el paracaídas mientras
ese valor pese menos de \SI{24.5}{\percent} del total, y el peso base
(\SI{24.3}{\percent}) queda apenas por debajo. Con los pesos de la tabla~\ref{tab:pugh}
el AHP elige el paracaídas con claridad.

\begin{nota}[Alcance de los números]
No se hizo ningún ensayo. Las velocidades salen de modelos de bibliografía; las masas y
volúmenes son estimaciones de diseño preliminar; las probabilidades de falla son
estimaciones de ingeniería por vuelo, coherentes con el capítulo~\ref{cap:confiabilidad}.
Los juicios del AHP sobre video y valor técnico son opiniones del equipo hechas
explícitas. El valor del capítulo está en mostrar \emph{qué supuestos deciden} la
elección, no en el tercer decimal de una prioridad.
\end{nota}

%=====================================================================
\section{Presupuesto de resistencia y envolvente}\label{sec:decision-presupuesto}
Toda alternativa debe aportar $C_DA\approx\SI{0.275}{m^2}$ para bajar a \SI{5}{m/s}. En
descenso estacionario el peso iguala la resistencia de todo el conjunto:
\begin{equation}
  W=\tfrac12\rho V^2\Big[(C_DA)_2+f\,C_{D,d}A_d+C_{D,c}A_c\Big]
  \quad\Rightarrow\quad
  V=\sqrt{\frac{2W}{\rho\,\big[(C_DA)_2+f\,C_{D,d}A_d+C_{D,c}A_c\big]}},
  \label{eq:decision-vterm}
\end{equation}
donde $(C_DA)_2$ es el aporte del sistema de la etapa~2 y $f\in[0,1]$ es la fracción del
aporte del drogue, que sigue atado y puede quedar en la estela. Con $W=\SI{4.905}{N}$ y
$\rho=\SI{1.225}{kg/m^3}$, la resistencia total para \SI{5}{m/s} es
$C_DA=2W/(\rho V^2)=\SI{0.320}{m^2}$. El drogue aporta \SI{0.0393}{m^2} y el cuerpo
\SI{0.0057}{m^2}; quedan \SI{0.275}{m^2} para la etapa~2. Para el límite superior de
\Req{C4} (\SI{8}{m/s}) bastan \SI{0.080}{m^2}.

Las alternativas se comparan en la misma envolvente: cuerpo Ø\,\SI{95}{mm} ×
\SI{250}{mm} (\Req{S2}, \Req{S3}), masa total \SI{500}{g} (\Req{S1}), sin pirotecnia
(\Req{M1}) y con el drogue atado. A los elementos \emph{rígidos} desplegables (palas,
aletas) se les aplica el mismo radio máximo que al rotor, $R\le\SI{0.20}{m}$. Los
elementos de tela no tienen ese límite, porque se pliegan dentro del cuerpo. Esta
asimetría es la clave de todo el capítulo: una tela puede ser tan grande como haga falta,
un elemento rígido no.

\paragraph{Incertidumbre.} Cada velocidad se propaga por Monte Carlo
($2\times10^5$ muestras) con distribuciones uniformes comunes: masa
\SIrange{490}{510}{g}, densidad \SIrange{1.10}{1.25}{kg/m^3} (casi todo el rango de \SIrange{0}{700}{m} y
\SIrange{5}{30}{\celsius}, que va de \SIrange{1.07}{1.27}{kg/m^3}), $C_{D,d}$ entre $0{,}75$ y $0{,}85$, y $f$ entre 0 y 1. A eso
se suma la incertidumbre propia de cada alternativa (tabla~\ref{tab:decision-cuant}).

\begin{figure}[H]
\centering
\begin{tikzpicture}
\begin{axis}[width=0.86\textwidth, height=6.4cm, xmin=0.05, xmax=0.45, ymin=2, ymax=14,
  xlabel={Radio desplegado característico $r$ [m]}, ylabel={$V$ estacionaria [m/s]},
  xtick={0.05,0.1,...,0.45}, xticklabel style={/pgf/number format/fixed, /pgf/number format/precision=2},
  grid=both, grid style={black!8}, legend pos=north east, legend cell align=left,
  legend style={font=\scriptsize, fill=white, fill opacity=0.9, text opacity=1}]
  \fill[ccuerda!12] (axis cs:0.05,2) rectangle (axis cs:0.45,8);
  \draw[ccuerda, dashed] (axis cs:0.05,5) -- (axis cs:0.45,5);
  \draw[cgris, densely dotted, line width=0.8pt] (axis cs:0.2,2) -- (axis cs:0.2,14)
    node[pos=0.93, left, font=\scriptsize, text=cgris] {$R=\SI{0.20}{m}$};
  \addplot[cfijo, thick] table[x=r, y=Vpar] {analysis/data/decision_tamano.dat};
  \addlegendentry{(a) paracaídas, radio nominal}
  \addplot[crot, thick] table[x=r, y=Vrot] {analysis/data/decision_tamano.dat};
  \addlegendentry{(c,e) rotor, $\CR=1{,}2$}
  \addplot[crot, densely dashed, thick] table[x=r, y=Vmono] {analysis/data/decision_tamano.dat};
  \addlegendentry{(d) monocóptero, $\CR=0{,}9$}
  \addplot[cfreno, thick] table[x=r, y=Val] {analysis/data/decision_tamano.dat};
  \addlegendentry{(f) 4 aletas de \SI{65}{mm} de ancho}
  \addplot[ccont, thick] table[x=r, y=Vesf] {analysis/data/decision_tamano.dat};
  \addlegendentry{(g) esfera, $C_D=0{,}47$}
  \node[font=\scriptsize, text=ccuerda!70!black, anchor=south west] at (axis cs:0.33,2.1) {banda de \Req{C4}};
  \node[circle, fill=cfijo, inner sep=1.4pt] at (axis cs:0.336,5.0) {};
  \node[circle, fill=crot, inner sep=1.4pt] at (axis cs:0.2,6.40) {};
  \node[circle, fill=cfreno, inner sep=1.4pt] at (axis cs:0.2,9.40) {};
  \node[circle, fill=ccont, inner sep=1.4pt] at (axis cs:0.127,10.79) {};
\end{axis}
\end{tikzpicture}
\caption{Velocidad de descenso frente al tamaño radial desplegado, con el drogue
aportando por completo ($f=1$). Los puntos marcan el diseño de cada alternativa: el
paracaídas de Ø\,\SI{0.67}{m} nominal, el rotor de $R=\SI{0.20}{m}$, las aletas de
\SI{152}{mm} y la esfera que infla un cartucho de \SI{16}{g} de CO$_2$. Para el
paracaídas, $r$ es el radio nominal (tela extendida); inflado, su radio proyectado es
unos dos tercios de ese valor. Datos:
\texttt{decision\_alternativas.py}.}
\label{fig:decision-tamano}
\end{figure}

La figura~\ref{fig:decision-tamano} resume la física. Por unidad de radio
\emph{nominal}, el rotor ($\CR\approx1{,}2$ sobre $\pi R^2$) supera al paracaídas plano
($C_{D0}=0{,}75$--$0{,}80$ sobre la tela extendida) y a la esfera ($0{,}47$). La
comparación justa es otra: el paracaídas inflado tiene un diámetro proyectado de unos
$0{,}66$--$0{,}70\,D_0$~\cite{knacke1992}, y referido a esa área su coeficiente es
$C_{D0}(D_0/D_p)^2\approx1{,}6$--$1{,}8$, mayor que el del rotor. El rotor no es, entonces, el
dispositivo más eficiente por área ocupada. Su desventaja decisiva es otra: con $R$
limitado a \SI{0.20}{m} solo llega a \SI{6.4}{m/s}, mientras que una tela de radio
nominal \SI{0.34}{m} (unos \SI{0.23}{m} inflada) llega a \SI{5}{m/s}.

%=====================================================================
\section{Modelo cuantitativo de cada alternativa}\label{sec:decision-modelos}

\subsection{(a) Paracaídas}\label{sec:decision-par}
Un paracaídas plano de Ø\,\SI{0.67}{m} da \SI{5.3}{m/s} (mediana del Monte Carlo) con
unos \SI{50}{g}. Para un paracaídas plano circular o hexagonal, Knacke da
$C_{D0}\approx0{,}75$--$0{,}80$ referido al área nominal $S_0$~\cite{knacke1992}. Con
$C_{D0}=0{,}775$:
\begin{equation}
  S_0=\frac{(C_DA)_2}{C_{D0}}=\SI{0.355}{m^2},\qquad D_0=\sqrt{4S_0/\pi}=\SI{0.672}{m}.
  \label{eq:decision-par}
\end{equation}
La masa se estima con tela de nailon \emph{ripstop} de \SI{40}{g/m^2} más
\SI{30}{\percent} por costuras (\SI{18}{g}), ocho líneas de $\approx D_0$ a
\SI{0.6}{g/m} (\SI{3}{g}), destorcedor (\SI{4}{g}), traba con servo (\SI{16}{g}) y
alojamiento (\SI{8}{g}): unos \SI{50}{g}. Con una densidad de plegado a mano de
\SI{0.35}{g/cm^3}, la tela y las líneas ocupan \SI{62}{cm^3}; con el alojamiento,
$\approx\SI{92}{cm^3}$ (un cilindro de Ø\,\SI{80}{mm} y \SI{18}{mm} de alto).

\paragraph{Oscilación.} El paracaídas plano es el tipo que más oscila: Knacke da
amplitudes del orden de $\pm\ang{10}$ a $\pm\ang{40}$, frente a $\pm\ang{10}$--$\ang{15}$
del hemisférico y amplitudes pequeñas de la cruz o la superficie guía~\cite{knacke1992}.
Son datos de paracaídas grandes; a esta escala pueden diferir. Como un
péndulo de largo efectivo $L\approx\SI{1}{m}$, el período es $T=2\pi\sqrt{L/g}\approx
\SI{2}{s}$. Con $\theta_0=\pm\ang{20}$, la velocidad angular de pico es
$\dot\theta_{\max}=2\pi\theta_0/T\approx\SI{60}{\degree/s}$: el video de nadir
(\Req{C7}) barre el suelo con un vaivén lento pero visible.

\paragraph{Riesgo propio.} Dos telas en la misma cuerda. La principal sale de su
alojamiento a \SI{13.4}{m/s} con el drogue arriba, y puede enredarse con él o con sus
líneas. La carga de apertura es baja ($q=\SI{110}{Pa}$; con $C_DS=\SI{0.275}{m^2}$ y un
factor de apertura de $1{,}5$, $F\approx\SI{45}{N}$). Se estima una probabilidad de falla
funcional del \SI{3}{\percent} por vuelo, dominada por el enredo.

\subsection{(b) Parapente pequeño}\label{sec:decision-pp}
Un ala inflada por presión dinámica (\emph{ram-air}) baja a \SI{5}{m/s} con
\SI{0.15}{m^2}, pero deriva unos \SI{720}{m} sin viento. En planeo estacionario la fuerza
aerodinámica total equilibra el peso. Con el drogue y el cuerpo arrastrados por la misma
línea, la resistencia es $D=\tfrac12\rho V^2(C_DS+k)$, con $k=C_{D,d}A_d+C_{D,c}A_c$, y
\begin{equation}
  V=\sqrt{\frac{2W}{\rho\sqrt{(C_LS)^2+(C_DS+k)^2}}},\qquad
  \frac{V_z}{V}=\frac{C_DS+k}{\sqrt{(C_LS)^2+(C_DS+k)^2}},\qquad
  \frac{V_x}{V_z}=\frac{C_LS}{C_DS+k}.
  \label{eq:decision-pp}
\end{equation}
Con $C_L=0{,}65$ y $C_D=0{,}20$ (ala chica de baja relación de aspecto, con
líneas~\cite{decision-lingard1995}), $V_z=\SI{5}{m/s}$ exige $S=\SI{0.147}{m^2}$. El
drogue arrastrado baja la fineza efectiva de $3{,}2$ a $1{,}28$: la velocidad horizontal
propia es \SI{6.4}{m/s}. Desde una liberación a \SI{560}{m} (80\,\% de un apogeo de
\SI{700}{m}) eso son \SI{720}{m} de deriva, sumados a los del viento. Sin el drogue
bastaría $S=\SI{0.071}{m^2}$, pero la deriva crece a \SI{1300}{m}.

El ala tiene unos \SI{48}{g} (tela con celdas $\approx2{,}6\,S$, 16 líneas) y ocupa
\SI{97}{cm^3} plegada. A esta escala la apertura de un ala con celdas es poco
confiable, y el rumbo no se controla: la deriva puede sacar el CanSat del campo. Se estima
una falla funcional del \SI{12}{\percent}.

\subsection{(c) Autogiro de cuatro palas (diseño base)}\label{sec:decision-rot}
El autogiro da \SI{6.4}{m/s} nominal y \SI{7.0}{m/s} de mediana (\SIrange{6.5}{7.7}{m/s}
al \SI{90}{\percent}), con \SI{70}{g}. La velocidad sale de la
ecuación~\eqref{eq:equilibrio} con $\CR$ uniforme entre $1{,}0$ y $1{,}3$ y
$A_R=\SI{0.1257}{m^2}$ (capítulo~\ref{cap:teoria}). La mediana del Monte Carlo es mayor
que el valor nominal porque incluye densidades menores que la del nivel del mar y un
drogue que puede no aportar ($f<1$). La masa de la etapa~2 es la de la
tabla~\ref{tab:masa}: palas (\SI{30}{g}), cubo, eje y rodamientos (\SI{24}{g}) y traba
(\SI{16}{g}). Las palas cuelgan por fuera del cuerpo, así que el volumen interno es
pequeño (\SI{40}{cm^3}, cubo y copa). Se cuentan siete piezas móviles: cuatro bisagras, el
cubo sobre su par de rodamientos, el anillo de traba y el destorcedor.

El cuerpo no gira con el rotor, pero el par de fricción de los rodamientos lo hace girar
despacio en guiñada. Sin datos, se estima una velocidad angular de unos
\SI{25}{\degree/s}, que la IMU medirá en los ensayos (\Req{SN5}). La probabilidad de falla
funcional (no arranca, traba, palas) se toma en \SI{4}{\percent}, coherente con el
\SI{3.3}{\percent} de no cumplir \Req{C4} del capítulo~\ref{cap:confiabilidad}.

\subsection{(d) Monocóptero tipo sámara}\label{sec:decision-mono}
Una sola pala da \SI{7.1}{m/s} nominal, pero hace girar todo el CanSat a
\SIrange{11}{23}{Hz}. Norberg~\cite{norberg1973} y Azuma y Yasuda~\cite{azuma1989}
muestran que la sámara se comporta como un disco de área $\pi R^2$ con una carga de disco
muy baja. Aquí el disco es el mismo ($R=\SI{0.20}{m}$), pero con una sola pala la solidez
es baja y el centro de giro no coincide con el eje del cuerpo; se toma $\CR$ entre
$0{,}7$ y $1{,}1$ (nominal $0{,}9$). La mediana es \SI{7.8}{m/s} y solo el
\SI{65}{\percent} de los casos queda bajo \SI{8}{m/s}.

Con una relación de velocidad de punta $\lambda=\Omega R/V$ entre 2 y 4, típica de
semillas~\cite{azuma1989}, el cuerpo gira a $f=\lambda V/(2\pi R)$, es decir, \SIrange{11}{23}{Hz}
(\SI{4000}{\degree/s} a \SI{8000}{\degree/s}). Ninguna cámara da un video útil a esas
velocidades, y el drogue necesita un destorcedor que trabaje de forma continua. Es la
alternativa más liviana (\SI{42}{g}) y una de las más simples.

\subsection{(e) Rotores coaxiales contrarrotantes}\label{sec:decision-coax}
Dos rotores en el mismo eje dan prácticamente la misma velocidad que uno
(\SI{6.3}{m/s}) con el doble de rotor (\SI{134}{g}). La teoría de cantidad de
movimiento explica por qué: los dos rotores actúan sobre el \emph{mismo} tubo de
corriente, y en esa teoría la resistencia del disco actuador depende de su área y de la
velocidad inducida, no del número de palas. El rotor inferior trabaja en la estela del superior y la interferencia aumenta
las pérdidas~\cite{decision-coleman1997,leishman2006}. Se toma $\CR$ entre $1{,}0$ y
$1{,}35$: algo más de solidez, sin ganancia de área. La ventaja real es otra: el par neto
sobre el cuerpo es casi nulo, y el video queda más estable que con un solo rotor.

El costo es alto: 12 piezas móviles, dos cubos apilados que consumen unos
\SI{30}{mm} del largo (\Req{S3}) y ocho palas de \SI{45}{mm} de cuerda alrededor de un
perímetro de \SI{276}{mm}, que no caben en una sola capa ($8\times45=\SI{360}{mm}$). Se
estima una falla funcional del \SI{8}{\percent}.

\subsection{(f) Aletas o cintas}\label{sec:decision-aletas}
Cuatro aletas rígidas dentro de $R\le\SI{0.20}{m}$ solo bajan a \SI{9.4}{m/s}: no cumplen
\Req{C4}. Cada aleta puede medir como máximo \SI{152}{mm} de largo radial y \SI{65}{mm} de
ancho, para que cuatro quepan plegadas en el perímetro. Abiertas a \ang{90} del flujo,
con $C_D\approx1{,}15$ (placa normal con interferencia del cuerpo~\cite{hoerner1965}),
aportan $C_DA=\SI{0.046}{m^2}$. Para \SI{5}{m/s} harían falta aletas de
\SI{0.92}{m} de largo. El Monte Carlo da \SIrange{9.6}{12.7}{m/s} y ningún caso bajo
\SI{8}{m/s}. Con cintas el resultado es peor: su $C_D$ referido al área de tela es
mucho menor que el de una placa normal, y harían falta varios metros cuadrados de tela
flameando junto al drogue.

\subsection{(g) Cuerpo inflable}\label{sec:decision-inf}
Una esfera que dé \SI{5}{m/s} necesitaría \SI{631}{g} de CO$_2$: más que todo el CanSat.
Con $C_D=0{,}47$ y $(C_DA)_2=\SI{0.275}{m^2}$,
\begin{equation}
  D=\sqrt{\frac{4(C_DA)_2}{\pi C_D}}=\SI{0.864}{m},\qquad
  \mathcal{V}=\frac{\pi D^3}{6}=\SI{337}{L},\qquad
  m_{\mathrm{CO_2}}=\rho_{\mathrm{CO_2}}\mathcal{V}=\SI{631}{g}.
  \label{eq:decision-esfera}
\end{equation}
Lo que cabe es un cartucho de \SI{16}{g} (\SI{8.6}{L} a presión ambiente): una esfera de
Ø\,\SI{0.25}{m} que da \SI{10.8}{m/s}, con unos \SI{98}{g} de cartucho, envoltura,
válvula y traba. Además, a \SI{5}{m/s} la esfera grande tendría
$Re=VD/\nu\approx\num{2.9e5}$, cerca de la crisis de resistencia: su $C_D$ podría caer a
la mitad. Un inflable por presión dinámica (\emph{ballute}) no tiene este problema, pero
entonces es un paracaídas con otra forma. Un generador de gas químico está prohibido por
\Req{M1}.

\subsection{Resumen cuantitativo}
\begin{table}[H]
\centering\small
\caption{Modelo cuantitativo de las alternativas dentro de la envolvente del CanSat.
$V$: mediana y percentiles 5--95 del Monte Carlo; $P_{\mathrm{in}}=P(\SI{2}{}\le V\le
\SI{8}{m/s})$ con el sistema funcionando; $P_F$: probabilidad de falla funcional por
vuelo (estimación); $\dot\theta$: movimiento angular típico del cuerpo (estimación).}
\label{tab:decision-cuant}
\setlength{\tabcolsep}{3.5pt}
\begin{tabular}{@{}llcccccccc@{}}
\toprule
 & \textbf{Tamaño} & \textbf{$V$} [m/s] & \textbf{$P_{\mathrm{in}}$} & \textbf{$|V-5|$} &
 \textbf{Masa} [g] & \textbf{Vol.} [cm$^3$] & \textbf{Piezas} & \textbf{$P_F$} & \textbf{$\dot\theta$} [°/s]\\\midrule
(a) & Ø\,\SI{0.67}{m} & \num{5.3} (\num{5.0}--\num{5.6}) & \num{1.00} & \num{0.29} & 50 & 92 & 2 & \num{0.03} & 60\\
(b) & \SI{0.147}{m^2} & \num{4.3} (\num{2.9}--\num{5.8}) & \num{1.00} & \num{0.92} & 48 & 97 & 2 & \num{0.12} & 20\\
(c) & $R=\SI{0.20}{m}$ & \num{7.0} (\num{6.5}--\num{7.7}) & \num{1.00} & \num{2.04} & 70 & 40 & 7 & \num{0.04} & 25\\
(d) & $R=\SI{0.20}{m}$ & \num{7.8} (\num{7.0}--\num{8.8}) & \num{0.65} & \num{2.82} & 42 & 30 & 4 & \num{0.15} & $>\num{4000}$\\
(e) & $2\times R=\SI{0.20}{m}$ & \num{7.0} (\num{6.4}--\num{7.6}) & \num{1.00} & \num{1.98} & 134 & 70 & 12 & \num{0.08} & 10\\
(f) & $4\times\SI{152}{mm}$ & \num{10.9} (\num{9.6}--\num{12.7}) & \num{0.00} & \num{5.98} & 60 & 10 & 5 & \num{0.05} & 30\\
(g) & Ø\,\SI{0.25}{m} & \num{13.3} (\num{11.3}--\num{17.0}) & \num{0.00} & \num{8.67} & 98 & 60 & 3 & \num{0.20} & 40\\
\bottomrule
\end{tabular}
\end{table}

\begin{figure}[H]
\centering
\begin{tikzpicture}
\begin{axis}[width=0.86\textwidth, height=5.6cm, ymin=2, ymax=18, xmin=-0.6, xmax=6.6,
  xtick={0,...,6}, xticklabels={(a) paracaídas,(b) parapente,(c) autogiro,(d) monocóptero,
  (e) coaxial,(f) aletas,(g) inflable},
  xticklabel style={font=\scriptsize, rotate=20, anchor=north east},
  ylabel={$V$ [m/s]}, ymajorgrids, grid style={black!8}]
  \fill[ccuerda!12] (axis cs:-0.6,2) rectangle (axis cs:6.6,8);
  \draw[ccuerda, dashed] (axis cs:-0.6,5) -- (axis cs:6.6,5);
  \addplot[only marks, mark=*, mark size=2.2pt, color=cfijo,
    error bars/.cd, y dir=both, y explicit, error bar style={line width=1pt, cfijo}]
    table[x=idx, y=Vmed, y error plus=errhi, y error minus=errlo]
    {analysis/data/decision_alternativas.dat};
\end{axis}
\end{tikzpicture}
\caption{Velocidad de la etapa~2 con el sistema funcionando: mediana e intervalo del
\SI{90}{\percent} (Monte Carlo). La banda verde es \Req{C4}.}
\label{fig:decision-vmc}
\end{figure}

\begin{resultado}[Filtro de factibilidad]
Las aletas (f) y el inflable (g) no cumplen \Req{C4} dentro de la envolvente: dan
\SI{10.9}{m/s} y \SI{13.3}{m/s}. El monocóptero (d) queda en el borde
($P_{\mathrm{in}}=0{,}65$) y anula el video. Las tres se mantienen en el AHP para
mostrar su efecto en el método, pero ninguna es candidata real. La decisión verdadera es
entre (a) paracaídas, (c) autogiro y, más atrás, (e) coaxial y (b) parapente.
\end{resultado}

%=====================================================================
\section{Proceso analítico jerárquico (AHP)}\label{sec:decision-ahp}

\subsection{Método}
El AHP de Saaty~\cite{decision-saaty1980,decision-saaty1990} ordena el problema en tres
niveles (figura~\ref{fig:decision-jerarquia}) y convierte juicios de a pares en
prioridades. Para $n$ elementos se arma una matriz recíproca $\mathbf{A}=[a_{ij}]$ con
$a_{ji}=1/a_{ij}$, donde $a_{ij}$ dice cuántas veces es preferible $i$ a $j$ en la escala
1--9 (1: igual; 3: moderada; 5: fuerte; 7: muy fuerte; 9: extrema). Las prioridades son
el vector propio principal:
\begin{equation}
  \mathbf{A}\,\mathbf{w}=\lambda_{\max}\mathbf{w},\qquad \sum_i w_i=1 .
  \label{eq:decision-eig}
\end{equation}
Si los juicios fueran perfectamente coherentes ($a_{ik}=a_{ij}a_{jk}$), sería
$\lambda_{\max}=n$. El índice y la razón de consistencia miden el apartamiento:
\begin{equation}
  CI=\frac{\lambda_{\max}-n}{n-1},\qquad CR=\frac{CI}{RI_n},
  \label{eq:decision-cr}
\end{equation}
con el índice aleatorio de Saaty $RI_n$ ($RI_7=1{,}32$). Se acepta $CR<0{,}10$. La
prioridad global de la alternativa $i$ es la suma ponderada (modo distributivo):
\begin{equation}
  G_i=\sum_{k=1}^{7} W_k\,p_{ik},
  \label{eq:decision-global}
\end{equation}
donde $W_k$ es el peso del criterio $k$ y $p_{ik}$ la prioridad local de $i$ en $k$.

\begin{figure}[H]
\centering
\begin{tikzpicture}[font=\scriptsize,
  nivel/.style={draw, rounded corners=2pt, align=center, inner sep=2.5pt, fill=white,
    minimum height=10.5mm, text width=15.5mm},
  lin/.style={draw=black!40, line width=0.35pt}]
  \node[nivel, draw=cfijo, fill=ffijo, text width=60mm, line width=0.8pt] (meta) at (0,3.0)
    {\textbf{Objetivo}: elegir el sistema de la etapa~2\\(\Req{C4} con drogue atado)};
  \foreach \k/\t [count=\i from 0] in {K1/Cumplir\\\Req{C4},K2/Riesgo\\funcional,K3/Masa,
     K4/Volumen\\interno,K5/Comple-\\jidad,K6/Video y\\deriva,K7/Valor\\técnico}{
    \node[nivel, draw=crot, fill=frot] (c\i) at ({(\i-3)*2.05},1.2) {\textbf{\k}\\\t};
    \draw[lin] (meta.south) -- (c\i.north);}
  \foreach \a/\t [count=\i from 0] in {a/Paracaídas,b/Parapente,c/Autogiro,d/Mono-\\cóptero,
     e/Coaxial,f/Aletas,g/Inflable}{
    \node[nivel, draw=ccuerda, fill=ccuerda!10] (a\i) at ({(\i-3)*2.05},-1.0) {(\a)\\\t};}
  \draw[lin] ($(c0.south)+(0,-0.15)$) -- ($(c6.south)+(0,-0.15)$);
  \draw[lin] ($(a0.north)+(0,0.15)$) -- ($(a6.north)+(0,0.15)$);
  \foreach \i in {0,...,6}{
    \draw[lin] (c\i.south) -- ++(0,-0.15);
    \draw[lin] (a\i.north) -- ++(0,0.15);}
  \draw[lin] ($(c3.south)+(0,-0.15)$) -- ($(a3.north)+(0,0.15)$);
  \node[font=\tiny, text=cgris, fill=white, inner sep=1pt] at (0,0.1) {cada criterio compara las 7 alternativas};
\end{tikzpicture}
\caption{Jerarquía del AHP: objetivo, siete criterios y siete alternativas.}
\label{fig:decision-jerarquia}
\end{figure}

\subsection{Criterios y cómo se obtienen los juicios}
Cinco de los siete criterios tienen una métrica del modelo; dos son juicios del equipo
(tabla~\ref{tab:decision-criterios}). Para los criterios medibles, la métrica $u_i$ se
lleva a una escala 1--9 y la diferencia de puntajes se lee como intensidad de preferencia:
\begin{equation}
  s_i=1+8\,\frac{u_i-u_{\min}}{u_{\max}-u_{\min}},\qquad
  a_{ij}=\mathcal{S}\big(1+s_i-s_j\big)\ \text{si } s_i\ge s_j,
  \label{eq:decision-escala}
\end{equation}
donde $\mathcal{S}(\cdot)$ redondea al valor más cercano de la escala de Saaty. Una
diferencia de 2 puntos es ``moderada'' y una de 8, ``extrema''. Para los
criterios de costo se usa $-u$. La métrica de \Req{C4} combina la probabilidad de
cumplir y la distancia al objetivo:
\begin{equation}
  u_1=P_{\mathrm{in}}\cdot\max\!\Big(0{,}02;\;1-\frac{\overline{|V-5|}}{\SI{6}{m/s}}\Big).
  \label{eq:decision-u1}
\end{equation}
El riesgo se mide en escala logarítmica, $u_2=-\log_{10}P_F$, porque pasar de
\SI{3}{\percent} a \SI{6}{\percent} de falla es tan grave como pasar de \SI{6}{\percent}
a \SI{12}{\percent}.

\begin{table}[H]
\centering\small
\caption{Criterios de decisión, métrica y origen de los juicios.}
\label{tab:decision-criterios}
\begin{tabularx}{\textwidth}{@{}clYl@{}}
\toprule
& \textbf{Criterio} & \textbf{Métrica o base del juicio} & \textbf{Origen}\\\midrule
K1 & Cumplir \Req{C4} & $u_1$, ecuación~\eqref{eq:decision-u1} & modelo\\
K2 & Riesgo funcional & $-\log_{10}P_F$ & estimación\\
K3 & Masa & masa del subsistema de etapa~2 [g] (costo) & modelo\\
K4 & Volumen interno & volumen plegado dentro del cuerpo [cm$^3$] (costo) & modelo\\
K5 & Complejidad & piezas móviles (costo) & diseño\\
K6 & Video y deriva & movimiento angular del cuerpo (\Req{C6}, \Req{C7}) y deriva horizontal & juicio\\
K7 & Valor técnico & contenido de ingeniería, aprendizaje, afinidad con las misiones de autogiro~\cite{cansat2019,cansat2025} y objetivo declarado del equipo & juicio\\
\bottomrule
\end{tabularx}
\end{table}

Para K6 y K7 el equipo asigna un puntaje ordinal de 1 a 9 y se aplica la misma
regla~\eqref{eq:decision-escala}. En K6: (e) 7, porque el par neto es nulo; (c) 6;
(b) 4, estable pero con deriva grande; (f) 4; (a) 3, por la oscilación pendular;
(g) 3, por la estela oscilante de la esfera; (d) 1, porque gira a más de \SI{10}{Hz}. En
K7: (c) 7, (e) 6, (d) 5, (b) 4, (g) 3, (f) 2 y (a) 1. El coaxial no puntúa más que el
autogiro en K7: duplicar un mecanismo no agrega conocimiento.

\subsection{Pesos de los criterios}
La matriz de criterios (tabla~\ref{tab:decision-wcrit}) da a cumplir \Req{C4}, al
riesgo y al valor técnico la misma importancia, y los pone por encima del resto; masa y
volumen pesan poco, porque la configuración~A tiene \SI{83}{g} de lastre disponible. El
resultado es $CR=0{,}009$. Esta matriz \emph{no} reproduce los pesos de la
tabla~\ref{tab:pugh}: allí \Req{C4} vale \SI{25}{\percent}, el riesgo \SI{20}{\percent} y
el valor técnico \SI{15}{\percent}. El equipo sube aquí el valor técnico a
\SI{24}{\percent} de forma deliberada, y ese juicio es el que decide el resultado
(sección~\ref{sec:decision-sens}). Por eso se evalúa también un escenario con los
\emph{pesos de la Pugh} trasladados a los siete criterios (K1 0,25; K2 0,20; K3 y K4
0,075 cada uno; K5 0,15; K6 0,10, que toma el lugar de la compatibilidad con el drogue;
K7 0,15). Se definen además otros dos escenarios: \emph{seguridad
primero} (el valor técnico baja a \SI{6}{\percent}) y \emph{misión técnica} (sube a
\SI{30}{\percent}), y el caso de pesos iguales.

\begin{table}[H]
\centering\small
\caption{Matriz de comparación de criterios (escenario base), pesos resultantes y pesos de
los escenarios alternativos. $\lambda_{\max}=7{,}072$, $CR=0{,}009$ (seguridad:
$0{,}011$; misión: $0{,}014$).}
\label{tab:decision-wcrit}
\setlength{\tabcolsep}{4pt}
\begin{tabular}{@{}l*{7}{c}cccc@{}}
\toprule
 & K1 & K2 & K3 & K4 & K5 & K6 & K7 & $W$ base & Seguridad & Misión & Iguales\\\midrule
K1 & 1 & 1 & 3 & 4 & 3 & 3 & 1 & \num{0.234} & \num{0.287} & \num{0.201} & \num{0.143}\\
K2 & 1 & 1 & 3 & 4 & 2 & 3 & 1 & \num{0.221} & \num{0.286} & \num{0.187} & \num{0.143}\\
K3 & 1/3 & 1/3 & 1 & 1 & 1/2 & 1/2 & 1/4 & \num{0.059} & \num{0.071} & \num{0.059} & \num{0.143}\\
K4 & 1/4 & 1/4 & 1 & 1 & 1/2 & 1/2 & 1/4 & \num{0.054} & \num{0.064} & \num{0.054} & \num{0.143}\\
K5 & 1/3 & 1/2 & 2 & 2 & 1 & 1 & 1/3 & \num{0.097} & \num{0.121} & \num{0.093} & \num{0.143}\\
K6 & 1/3 & 1/3 & 2 & 2 & 1 & 1 & 1/3 & \num{0.092} & \num{0.114} & \num{0.108} & \num{0.143}\\
K7 & 1 & 1 & 4 & 4 & 3 & 3 & 1 & \num{0.243} & \num{0.058} & \num{0.297} & \num{0.143}\\
\bottomrule
\end{tabular}
\end{table}

\subsection{Matrices por criterio}
Las siete matrices de alternativas son coherentes ($CR$ entre $0{,}020$ y $0{,}051$). La
tabla~\ref{tab:decision-k2} muestra como ejemplo la de riesgo, y la
tabla~\ref{tab:decision-locales}, las prioridades locales de todas. Las matrices
completas las imprime \texttt{decision\_ahp.py}.

\begin{table}[H]
\centering\small
\caption{Matriz de comparación por pares para K2 (riesgo funcional), derivada de $P_F$
con la ecuación~\eqref{eq:decision-escala}. $\lambda_{\max}=7{,}311$, $CR=0{,}039$.}
\label{tab:decision-k2}
\begin{tabular}{@{}l*{7}{c}c@{}}
\toprule
 & (a) & (b) & (c) & (d) & (e) & (f) & (g) & $p_{i2}$\\\midrule
(a) $P_F=0{,}03$ & 1 & 7 & 2 & 8 & 5 & 3 & 9 & \num{0.368}\\
(b) $0{,}12$ & 1/7 & 1 & 1/6 & 2 & 1/3 & 1/5 & 3 & \num{0.047}\\
(c) $0{,}04$ & 1/2 & 6 & 1 & 7 & 4 & 2 & 8 & \num{0.256}\\
(d) $0{,}15$ & 1/8 & 1/2 & 1/7 & 1 & 1/4 & 1/6 & 2 & \num{0.033}\\
(e) $0{,}08$ & 1/5 & 3 & 1/4 & 4 & 1 & 1/3 & 5 & \num{0.093}\\
(f) $0{,}05$ & 1/3 & 5 & 1/2 & 6 & 3 & 1 & 7 & \num{0.179}\\
(g) $0{,}20$ & 1/9 & 1/3 & 1/8 & 1/2 & 1/5 & 1/7 & 1 & \num{0.024}\\
\bottomrule
\end{tabular}
\end{table}

\begin{table}[H]
\centering\small
\caption{Prioridades locales $p_{ik}$ y razón de consistencia de cada matriz.}
\label{tab:decision-locales}
\begin{tabular}{@{}l*{7}{c}@{}}
\toprule
 & K1 & K2 & K3 & K4 & K5 & K6 & K7\\\midrule
(a) Paracaídas  & \textbf{0,355} & \textbf{0,368} & 0,203 & 0,030 & \textbf{0,269} & 0,066 & 0,031\\
(b) Parapente   & 0,262 & 0,047 & 0,203 & 0,024 & \textbf{0,269} & 0,109 & 0,104\\
(c) Autogiro    & 0,137 & 0,256 & 0,092 & 0,159 & 0,051 & 0,252 & \textbf{0,354}\\
(d) Monocóptero & 0,058 & 0,033 & \textbf{0,307} & 0,227 & 0,118 & 0,031 & 0,159\\
(e) Coaxial     & 0,141 & 0,093 & 0,020 & 0,058 & 0,019 & \textbf{0,366} & 0,240\\
(f) Aletas      & 0,023 & 0,179 & 0,129 & \textbf{0,421} & 0,093 & 0,109 & 0,045\\
(g) Inflable    & 0,023 & 0,024 & 0,044 & 0,080 & 0,181 & 0,066 & 0,068\\\midrule
$CR$            & 0,049 & 0,039 & 0,034 & 0,051 & 0,039 & 0,020 & 0,025\\
\bottomrule
\end{tabular}
\end{table}

\subsection{Resultado del AHP}
Con los pesos base, el paracaídas y el autogiro empatan: $G_a=0{,}218$ y $G_c=0{,}217$
(figura~\ref{fig:decision-global}). La diferencia, una milésima, es mucho menor que la
incertidumbre de cualquier juicio. El paracaídas gana en \Req{C4}, riesgo y
complejidad; el autogiro, en video y valor técnico. Siguen el coaxial ($0{,}152$) y el
parapente ($0{,}147$). Con \emph{seguridad primero} el paracaídas gana con claridad
($0{,}265$ contra $0{,}185$); con \emph{misión técnica} gana el autogiro ($0{,}227$
contra $0{,}195$); con pesos iguales vuelven a empatar ($0{,}189$ y $0{,}186$). Con los pesos de la Pugh
gana el paracaídas con margen ($0{,}231$ contra $0{,}190$).

\begin{figure}[H]
\centering
\begin{tikzpicture}
\begin{axis}[width=0.92\textwidth, height=6cm, ybar, bar width=3.6pt, ymin=0, ymax=0.3,
  xtick={0,...,6}, xticklabels={(a),(b),(c),(d),(e),(f),(g)}, enlarge x limits=0.08,
  ylabel={Prioridad global $G_i$}, ymajorgrids, grid style={black!8},
  legend style={font=\scriptsize, at={(0.5,1.03)}, anchor=south, legend columns=5,
    /tikz/every even column/.append style={column sep=6pt}},
  yticklabel style={/pgf/number format/fixed, /pgf/number format/precision=2}]
  \addplot[fill=cfijo, draw=cfijo] table[x=idx, y=ahp_base] {analysis/data/decision_global.dat};
  \addlegendentry{base}
  \addplot[fill=cfreno!70, draw=cfreno] table[x=idx, y=ahp_seguridad] {analysis/data/decision_global.dat};
  \addlegendentry{seguridad primero}
  \addplot[fill=crot!80, draw=crot] table[x=idx, y=ahp_mision] {analysis/data/decision_global.dat};
  \addlegendentry{misión técnica}
  \addplot[fill=cgris!50, draw=cgris] table[x=idx, y=ahp_iguales] {analysis/data/decision_global.dat};
  \addlegendentry{pesos iguales}
  \addplot[fill=ccuerda!60, draw=ccuerda] table[x=idx, y=ahp_pugh] {analysis/data/decision_global.dat};
  \addlegendentry{pesos Pugh}
\end{axis}
\end{tikzpicture}
\caption{Prioridades globales del AHP en cinco escenarios de pesos.}
\label{fig:decision-global}
\end{figure}

%=====================================================================
\section{TOPSIS como segundo método}\label{sec:decision-topsis}
TOPSIS elige el autogiro con margen ($0{,}79$ contra $0{,}66$ del coaxial, $0{,}63$ del
parapente y $0{,}61$ del paracaídas). El método de Hwang y Yoon~\cite{decision-hwang1981} trabaja con las
métricas crudas $x_{ik}$ de la tabla~\ref{tab:decision-cuant} (y los puntajes de K6 y
K7). Normaliza cada columna, pondera, y mide la distancia de cada alternativa a una
solución ideal $v^+$ y a una anti-ideal $v^-$:
\begin{gather}
  r_{ik}=\frac{x_{ik}}{\sqrt{\sum_j x_{jk}^2}},\qquad v_{ik}=W_k\,r_{ik},\\
  D_i^{\pm}=\sqrt{\sum_k\big(v_{ik}-v_k^{\pm}\big)^2},\qquad
  C_i=\frac{D_i^-}{D_i^++D_i^-}\in[0,1],
  \label{eq:decision-topsis}
\end{gather}
donde $v_k^+$ es el mejor valor de la columna (máximo en beneficios, mínimo en costos) y
$v_k^-$ el peor.

\begin{table}[H]
\centering\small
\caption{Coeficiente de cercanía $C_i$ de TOPSIS en los cinco escenarios.}
\label{tab:decision-topsis}
\begin{tabular}{@{}l*{7}{c}@{}}
\toprule
Escenario & (a) & (b) & (c) & (d) & (e) & (f) & (g)\\\midrule
Base              & 0,605 & 0,625 & \textbf{0,788} & 0,458 & 0,656 & 0,423 & 0,258\\
Seguridad primero & \textbf{0,812} & 0,655 & 0,754 & 0,390 & 0,617 & 0,480 & 0,236\\
Misión técnica    & 0,523 & 0,601 & \textbf{0,807} & 0,493 & 0,674 & 0,396 & 0,284\\
Pesos iguales     & 0,591 & 0,594 & \textbf{0,715} & 0,530 & 0,502 & 0,545 & 0,395\\
Pesos Pugh        & 0,700 & 0,674 & \textbf{0,726} & 0,464 & 0,567 & 0,445 & 0,319\\
\bottomrule
\end{tabular}
\end{table}

\paragraph{Por qué TOPSIS y AHP difieren.} TOPSIS castiga las debilidades extremas. El
paracaídas es el \emph{peor} de todos en valor técnico, y casi el peor en volumen
interno; esa distancia al ideal pesa mucho con $W_7=0{,}24$. El autogiro, en cambio,
no es el peor en nada y queda cerca del ideal en casi todo. Además, con $P_F$ en escala
lineal, la diferencia entre \SI{3}{\percent} y \SI{4}{\percent} es pequeña frente al
\SI{20}{\percent} del inflable, que fija el anti-ideal. El AHP con
la regla~\eqref{eq:decision-escala} es más compensatorio y lee esa misma diferencia como
``igual a moderada''. Ninguno de los dos es ``el correcto'': miden cosas
distintas, y por eso conviene usar ambos. Con los pesos de la Pugh, TOPSIS todavía
prefiere el autogiro, pero por poco ($0{,}726$ frente a $0{,}700$).

\begin{table}[H]
\centering\small
\caption{Robustez frente al método, con los pesos base. Se marca la alternativa ganadora.}
\label{tab:decision-variantes}
\begin{tabularx}{\textwidth}{@{}Ycccc@{}}
\toprule
\textbf{Variante} & \textbf{(a)} & \textbf{(c)} & \textbf{(e)} & \textbf{Gana}\\\midrule
AHP, juicios por diferencia (base) & \textbf{0,218} & 0,217 & 0,152 & (a), por 0,001\\
AHP, juicios por cociente $a_{ij}=\mathcal{S}(s_i/s_j)$ & 0,156 & \textbf{0,211} & 0,165 & (c)\\
AHP, modo ideal & \textbf{0,220} & 0,211 & 0,147 & (a)\\
AHP sin (f) ni (g), métricas reescaladas & \textbf{0,286} & 0,255 & 0,166 & (a)\\
TOPSIS, métricas crudas (base) & 0,605 & \textbf{0,788} & 0,656 & (c)\\
TOPSIS, riesgo en escala logarítmica & 0,560 & \textbf{0,758} & 0,631 & (c)\\
TOPSIS con los puntajes 1--9 del AHP & 0,572 & \textbf{0,794} & 0,627 & (c)\\
TOPSIS sin (f) ni (g) & 0,555 & \textbf{0,770} & 0,592 & (c)\\
\bottomrule
\end{tabularx}
\end{table}

La tabla~\ref{tab:decision-variantes} muestra que TOPSIS elige el autogiro en todas sus
variantes, y que el AHP lo hace cambiar según detalles de método. Con juicios por
cociente, el AHP comprime las diferencias entre alternativas buenas (una métrica de
\num{9} contra \num{6.5} da $a_{ij}=1$) y favorece al autogiro. Quitar las dos
alternativas no factibles cambia la escala de las métricas y amplía la ventaja del
paracaídas (de $0{,}001$ a $0{,}031$). El orden entre (a) y (c) no se invierte en este
caso, pero el efecto es el mismo que origina la conocida inversión de rangos del AHP
distributivo: la prioridad de una alternativa depende de cuáles otras entran en la
comparación~\cite{decision-belton1983}.

%=====================================================================
\section{Sensibilidad}\label{sec:decision-sens}

\subsection{Monte Carlo sobre los pesos y los juicios}
En la línea de la simulación de pesos de Butler, Jia y Dyer~\cite{decision-butler1997},
los pesos se muestrean de una distribución de Dirichlet:
\begin{equation}
  \mathbf{W}\sim\mathrm{Dir}(\kappa\,\mathbf{W}_0),\qquad
  \mathrm{E}[W_k]=W_{0,k},\qquad
  \mathrm{Var}[W_k]=\frac{W_{0,k}(1-W_{0,k})}{\kappa+1}.
  \label{eq:decision-dir}
\end{equation}
Con $\kappa=30$, un peso base de $0{,}24$ varía con desvío $\pm0{,}077$: equivale a que
cada integrante del equipo opine distinto dentro de lo razonable. La Dirichlet uniforme
($\kappa\mathbf{W}_0=\mathbf{1}$) recorre todo el simplex: representa no saber nada de
los pesos. En un tercer caso se perturban también los juicios: cada $a_{ij}$ se multiplica
por $e^{\varepsilon}$ con $\varepsilon$ normal de media 0 y desvío $0{,}35$ (algo menos de un
escalón de la escala) y cada $P_F$ por un factor lognormal con factor de error 2. Se usan
\num{40000} muestras (\num{8000} en el caso con juicios).

\begin{table}[H]
\centering\small
\caption{Probabilidad de que cada alternativa sea la mejor.}
\label{tab:decision-pmejor}
\begin{tabular}{@{}l*{7}{c}@{}}
\toprule
Caso & (a) & (b) & (c) & (d) & (e) & (f) & (g)\\\midrule
AHP, $\mathrm{Dir}(30\,\mathbf{W}_0)$           & \textbf{0,52} & 0,00 & 0,48 & 0,00 & 0,00 & 0,00 & 0,00\\
AHP, $\mathrm{Dir}(30\,\mathbf{W}_0)$ y juicios & 0,45 & 0,00 & \textbf{0,52} & 0,00 & 0,02 & 0,01 & 0,00\\
AHP, Dirichlet uniforme                         & \textbf{0,46} & 0,01 & 0,32 & 0,04 & 0,04 & 0,12 & 0,00\\
TOPSIS, $\mathrm{Dir}(30\,\mathbf{W}_0)$        & 0,05 & 0,04 & \textbf{0,90} & 0,00 & 0,00 & 0,00 & 0,00\\
TOPSIS, Dirichlet uniforme                      & 0,16 & 0,14 & \textbf{0,55} & 0,05 & 0,01 & 0,09 & 0,00\\
\bottomrule
\end{tabular}
\end{table}

\begin{figure}[H]
\centering
\begin{tikzpicture}
\begin{axis}[width=0.8\textwidth, height=5.2cm, ybar, bar width=7pt, ymin=0, ymax=0.62,
  xtick={1,...,7}, xlabel={Puesto en el AHP}, ylabel={Frecuencia}, ymajorgrids,
  grid style={black!8}, enlarge x limits=0.08,
  legend style={font=\scriptsize, at={(0.98,0.97)}, anchor=north east},
  yticklabel style={/pgf/number format/fixed, /pgf/number format/precision=1}]
  \addplot[fill=cfijo, draw=cfijo] table[x=rango, y=a] {analysis/data/decision_rangos.dat};
  \addlegendentry{(a) paracaídas}
  \addplot[fill=crot, draw=crot] table[x=rango, y=c] {analysis/data/decision_rangos.dat};
  \addlegendentry{(c) autogiro}
  \addplot[fill=ccont!60, draw=ccont] table[x=rango, y=e] {analysis/data/decision_rangos.dat};
  \addlegendentry{(e) coaxial}
\end{axis}
\end{tikzpicture}
\caption{Distribución del puesto de (a), (c) y (e) en el AHP con pesos
$\mathrm{Dir}(30\,\mathbf{W}_0)$. El autogiro queda primero o segundo en el
\SI{98}{\percent} de los casos.}
\label{fig:decision-rangos}
\end{figure}

Tres lecturas. Primera: con pesos razonables solo (a) y (c) pueden ganar (uno de los dos es el primero
en el \SI{99.98}{\percent} de los casos); el autogiro queda entre los dos primeros en el
\SI{98}{\percent} de los casos y el paracaídas en el \SI{94}{\percent}
(figura~\ref{fig:decision-rangos}).
Segunda: dentro del AHP la elección es una moneda al aire; perturbar los juicios la
inclina apenas hacia el autogiro. Tercera: si no se sabe nada de los pesos (Dirichlet
uniforme), el paracaídas es la opción más probable en el AHP, y las aletas ganan en el
\SI{12}{\percent} de los casos. Eso ocurre cuando gran parte del peso cae en el volumen (en esos casos
$W_4$ vale $0{,}39$ en promedio):
la Dirichlet uniforme ignora que (f) no cumple \Req{C4}. Es un recordatorio de que el
filtro de factibilidad debe ir \emph{antes} de cualquier ponderación.

\subsection{Barridos de un peso y cambio de ranking}
Variar un peso a la vez muestra que el empate del AHP está en un filo. Se busca el
cambio mínimo de un peso que invierte el orden, como proponen Triantaphyllou y
Sánchez~\cite{decision-triantaphyllou1997}. Cada barrido fija
$W_k=x$ y reescala los demás en proporción, $W_j=W_{0,j}(1-x)/(1-W_{0,k})$
(figura~\ref{fig:decision-barridos}). La tabla~\ref{tab:decision-cruces} da el peso donde
el paracaídas y el autogiro se cruzan.

\begin{figure}[H]
\centering
\pgfplotsset{barr/.style={width=\linewidth, height=5.4cm, xmin=0, xmax=0.7,
  ymin=0, ymax=0.42, grid=both, grid style={black!8}, ylabel={$G_i$},
  tick label style={font=\scriptsize}, label style={font=\small},
  yticklabel style={/pgf/number format/fixed, /pgf/number format/precision=1},
  xticklabel style={/pgf/number format/fixed, /pgf/number format/precision=1}}}
\begin{subfigure}{0.49\textwidth}\centering
\begin{tikzpicture}
\begin{axis}[barr, xlabel={$W_7$ (valor técnico)}]
  \addplot[cfijo, thick] table[x=w, y=a] {analysis/data/decision_barrido_K7.dat};
  \addplot[crot, thick] table[x=w, y=c] {analysis/data/decision_barrido_K7.dat};
  \addplot[ccont, thick, densely dashed] table[x=w, y=e] {analysis/data/decision_barrido_K7.dat};
  \addplot[ccuerda, thick, densely dotted] table[x=w, y=b] {analysis/data/decision_barrido_K7.dat};
  \draw[cgris, densely dotted] (axis cs:0.243,0) -- (axis cs:0.243,0.42);
\end{axis}
\end{tikzpicture}
\caption{Valor técnico (K7).}
\end{subfigure}\hfill
\begin{subfigure}{0.49\textwidth}\centering
\begin{tikzpicture}
\begin{axis}[barr, xlabel={$W_2$ (riesgo)}]
  \addplot[cfijo, thick] table[x=w, y=a] {analysis/data/decision_barrido_K2.dat};
  \addplot[crot, thick] table[x=w, y=c] {analysis/data/decision_barrido_K2.dat};
  \addplot[ccont, thick, densely dashed] table[x=w, y=e] {analysis/data/decision_barrido_K2.dat};
  \addplot[ccuerda, thick, densely dotted] table[x=w, y=b] {analysis/data/decision_barrido_K2.dat};
  \draw[cgris, densely dotted] (axis cs:0.221,0) -- (axis cs:0.221,0.42);
\end{axis}
\end{tikzpicture}
\caption{Riesgo funcional (K2).}
\end{subfigure}

\vspace{2mm}
\begin{subfigure}{0.49\textwidth}\centering
\begin{tikzpicture}
\begin{axis}[barr, xlabel={$W_1$ (cumplir \Req{C4})}]
  \addplot[cfijo, thick] table[x=w, y=a] {analysis/data/decision_barrido_K1.dat};
  \addplot[crot, thick] table[x=w, y=c] {analysis/data/decision_barrido_K1.dat};
  \addplot[ccont, thick, densely dashed] table[x=w, y=e] {analysis/data/decision_barrido_K1.dat};
  \addplot[ccuerda, thick, densely dotted] table[x=w, y=b] {analysis/data/decision_barrido_K1.dat};
  \draw[cgris, densely dotted] (axis cs:0.234,0) -- (axis cs:0.234,0.42);
\end{axis}
\end{tikzpicture}
\caption{Cumplimiento de \Req{C4} (K1).}
\end{subfigure}\hfill
\begin{subfigure}{0.49\textwidth}\centering
\begin{tikzpicture}
\begin{axis}[barr, xlabel={$W_6$ (video y deriva)}, legend style={font=\scriptsize,
  at={(0.02,0.98)}, anchor=north west}]
  \addplot[cfijo, thick] table[x=w, y=a] {analysis/data/decision_barrido_K6.dat};
  \addlegendentry{(a)}
  \addplot[crot, thick] table[x=w, y=c] {analysis/data/decision_barrido_K6.dat};
  \addlegendentry{(c)}
  \addplot[ccont, thick, densely dashed] table[x=w, y=e] {analysis/data/decision_barrido_K6.dat};
  \addlegendentry{(e)}
  \addplot[ccuerda, thick, densely dotted] table[x=w, y=b] {analysis/data/decision_barrido_K6.dat};
  \addlegendentry{(b)}
  \draw[cgris, densely dotted] (axis cs:0.092,0) -- (axis cs:0.092,0.42);
\end{axis}
\end{tikzpicture}
\caption{Video y deriva (K6).}
\end{subfigure}
\caption{Prioridad global del AHP al variar un peso (los demás se reescalan en
proporción). La línea punteada vertical marca el peso base. Azul: paracaídas; naranja:
autogiro; violeta: coaxial; verde: parapente.}
\label{fig:decision-barridos}
\end{figure}

\begin{table}[H]
\centering\small
\caption{Cambios de ganador al barrer un peso (los demás se reescalan en proporción). Se
indica el ganador en cada tramo y el peso donde cambia. Los cruces del AHP y el de K7 en
TOPSIS están interpolados; los demás de TOPSIS, con paso de $0{,}01$.}
\label{tab:decision-cruces}
\begin{tabular}{@{}lcll@{}}
\toprule
Criterio & $W_0$ & AHP & TOPSIS\\\midrule
K1 Cumplir \Req{C4}  & 0,234 & (c) $\xrightarrow{0{,}232}$ (a) & (c) $\xrightarrow{0{,}40}$ (b) $\xrightarrow{0{,}55}$ (a)\\
K2 Riesgo            & 0,221 & (c) $\xrightarrow{0{,}215}$ (a) & (c) $\xrightarrow{0{,}81}$ (a)\\
K3 Masa              & 0,059 & (c) $\xrightarrow{0{,}053}$ (a) $\xrightarrow{0{,}56}$ (d) & (c) $\xrightarrow{0{,}38}$ (b) $\xrightarrow{0{,}75}$ (d)\\
K4 Volumen           & 0,054 & (a) $\xrightarrow{0{,}060}$ (c) $\xrightarrow{0{,}34}$ (f) & (c) $\xrightarrow{0{,}44}$ (f)\\
K5 Complejidad       & 0,097 & (c) $\xrightarrow{0{,}093}$ (a) & (c) $\xrightarrow{0{,}23}$ (b)\\
K6 Video y deriva    & 0,092 & (a) $\xrightarrow{0{,}096}$ (c) $\xrightarrow{0{,}43}$ (e) & (c) $\xrightarrow{0{,}45}$ (e)\\
K7 Valor técnico     & 0,243 & (a) $\xrightarrow{0{,}245}$ (c) & (a) $\xrightarrow{0{,}11}$ (c)\\
\bottomrule
\end{tabular}
\end{table}

En el AHP, todos los cruces entre (a) y (c) están a menos de $0{,}007$ de su peso base.
Cualquier cambio pequeño de opinión invierte el resultado: no hay una elección robusta.
En TOPSIS, en cambio, el autogiro gana en un entorno amplio del peso base. Deja de ganar
si el valor técnico pesa menos de $0{,}11$ (gana el paracaídas), si la complejidad pesa
más de $0{,}23$ o \Req{C4} más de $0{,}40$ (gana el parapente, que TOPSIS ve como un
compromiso: simple, liviano y con buena velocidad vertical, sin penalizar bien su deriva),
o si el riesgo pesa más de $0{,}81$. Que aparezcan (d), (f) o (e) con pesos extremos en
masa, volumen o video solo muestra que cada una es la mejor en un único aspecto.

\subsection{¿Cuánto riesgo tolera la elección del autogiro?}
El autogiro gana mientras su probabilidad de falla quede bajo \SI{3.5}{\percent} (AHP
base), \SI{6.8}{\percent} (AHP misión técnica) o \SI{19}{\percent} (TOPSIS base). Con la
estimación actual de \SI{4}{\percent} ya está por encima del umbral del AHP base. La figura~\ref{fig:decision-umbral} muestra la ventaja del
autogiro sobre el paracaídas al variar solo $P_{F,c}$.
Este es el número que los ensayos deben fijar: si las caídas de prueba
(capítulo~\ref{cap:riesgos}) indican una tasa de falla de arranque mayor que una cada 15 lanzamientos
(\SI{6.8}{\percent}), el autogiro deja de ser defendible aun con los pesos de la misión
técnica.

\begin{figure}[H]
\centering
\begin{tikzpicture}
\begin{semilogxaxis}[width=0.8\textwidth, height=4.6cm, xmin=0.01, xmax=0.3,
  ymin=-0.15, ymax=0.22, xlabel={$P_{F}$ del autogiro}, ylabel={Ventaja de (c) sobre (a)},
  grid=both, grid style={black!8}, legend style={font=\scriptsize, at={(0.5,1.03)},
  anchor=south, legend columns=2, /tikz/every even column/.append style={column sep=6pt}},
  log basis x=10,
  xtick={0.01,0.02,0.05,0.1,0.2},
  xticklabel={\pgfmathparse{exp(\tick*ln(10))}\pgfmathprintnumber[fixed,precision=2]{\pgfmathresult}},
  yticklabel style={/pgf/number format/fixed, /pgf/number format/precision=2}]
  \draw[black!60] (axis cs:0.01,0) -- (axis cs:0.3,0);
  \draw[cgris, densely dotted] (axis cs:0.04,-0.15) -- (axis cs:0.04,0.22)
    node[pos=0.75, right, font=\scriptsize, text=cgris] {estimación};
  \addplot[cfijo, thick] table[x=pf, y=dAHPbase] {analysis/data/decision_umbralpf.dat};
  \addlegendentry{AHP base: $G_c-G_a$}
  \addplot[crot, thick] table[x=pf, y=dAHPmision] {analysis/data/decision_umbralpf.dat};
  \addlegendentry{AHP misión técnica: $G_c-G_a$}
  \addplot[ccont, thick, densely dashed] table[x=pf, y=dTOPbase] {analysis/data/decision_umbralpf.dat};
  \addlegendentry{TOPSIS base: $C_c-C_a$}
\end{semilogxaxis}
\end{tikzpicture}
\caption{Ventaja del autogiro sobre el paracaídas en función de su probabilidad de falla
funcional. El autogiro gana mientras la curva esté sobre cero. Las curvas del AHP son
escalonadas porque los juicios se redondean a la escala de Saaty.}
\label{fig:decision-umbral}
\end{figure}

\subsection{Frontera de decisión}
El plano de la figura~\ref{fig:decision-frontera} cruza los dos criterios en disputa:
riesgo y valor técnico. El AHP elige el autogiro si $W_7\gtrsim0{,}24$--$0{,}25$, casi
sin importar $W_2$. TOPSIS lo elige con valores menores de $W_7$, tanto más cuanto más
pesa el riesgo, porque al subir $W_2$ baja el peso de \Req{C4}, que es donde el paracaídas
saca ventaja.

\begin{figure}[H]
\centering
\begin{tikzpicture}
\begin{axis}[width=0.72\textwidth, height=6cm, xmin=0, xmax=0.6, ymin=0, ymax=0.45,
  xlabel={$W_2$ (riesgo)}, ylabel={$W_7$ (valor técnico)}, grid=both,
  grid style={black!8}, legend style={font=\scriptsize, at={(0.98,0.98)}, anchor=north east},
  xticklabel style={/pgf/number format/fixed, /pgf/number format/precision=1},
  yticklabel style={/pgf/number format/fixed, /pgf/number format/precision=1}]
  \addplot[name path=ahp, crot, thick] table[x=w2, y=w7ahp] {analysis/data/decision_frontera.dat};
  \addlegendentry{frontera AHP}
  \addplot[name path=top, crot, thick, densely dashed] table[x=w2, y=w7top] {analysis/data/decision_frontera.dat};
  \addlegendentry{frontera TOPSIS}
  \path[name path=piso] (axis cs:0,0) -- (axis cs:0.6,0);
  \path[name path=techo] (axis cs:0,0.45) -- (axis cs:0.6,0.45);
  \addplot[cfijo!15] fill between[of=piso and top];
  \addplot[crot!12] fill between[of=ahp and techo];
  \node[circle, fill=black, inner sep=1.6pt, label={[font=\scriptsize]right:base}] at (axis cs:0.221,0.243) {};
  \node[circle, fill=cfreno, inner sep=1.4pt, label={[font=\scriptsize]right:seguridad}] at (axis cs:0.286,0.058) {};
  \node[circle, fill=crot, inner sep=1.4pt, label={[font=\scriptsize]right:misión}] at (axis cs:0.187,0.297) {};
  \node[font=\scriptsize, text=cfijo] at (axis cs:0.45,0.03) {gana (a) en ambos};
  \node[font=\tiny, text=ccuerda!70!black, align=center] at (axis cs:0.055,0.13) {TOPSIS:\\(b)};
  \node[font=\scriptsize, text=crot!80!black] at (axis cs:0.13,0.40) {gana (c) en ambos};
\end{axis}
\end{tikzpicture}
\caption{Frontera paracaídas/autogiro en el plano riesgo--valor técnico (los otros cinco
pesos conservan sus proporciones base; los puntos de los escenarios, que no cumplen
esa condición, se ubican de forma aproximada). Entre las dos curvas los métodos discrepan. En el
AHP gana siempre (a) o (c); en TOPSIS, bajo su curva, gana (a), salvo en la esquina de
riesgo bajo ($W_2\lesssim0{,}15$, $W_7\approx0{,}05$--$0{,}2$), donde gana el parapente.}
\label{fig:decision-frontera}
\end{figure}

%=====================================================================
\section{Decisión y su justificación}\label{sec:decision-final}

\begin{resultado}[Qué dice el análisis]
\begin{enumerate}
  \item En prestaciones puras, \textbf{el paracaídas es mejor}: baja a
        \SI{5.3}{m/s} (el autogiro, a \SI{7.0}{m/s}, cerca del límite de \SI{8}{m/s}),
        es más liviano, más simple y algo más confiable.
  \item El autogiro es mejor en video (sin péndulo de \SI{60}{\degree/s}), en volumen
        interno y, sobre todo, en valor técnico.
  \item AHP y TOPSIS coinciden en que, con pesos razonables, solo (a) y (c) ganan con
        probabilidad apreciable, y discrepan en cuál. El AHP los da empatados; TOPSIS
        prefiere el autogiro con margen.
  \item La elección depende de un solo juicio: cuánto vale el contenido técnico. Con
        $W_7<0{,}245$ en el AHP, o $W_7<0{,}11$ en TOPSIS, conviene el paracaídas. El
        peso base ($0{,}243$) está en el filo, y con los pesos de la tabla~\ref{tab:pugh}
        ($W_7=0{,}15$) el AHP elige el paracaídas.
\end{enumerate}
\end{resultado}

\paragraph{Cuándo debería ganar el paracaídas.} En cuatro escenarios concretos: (1) si el
objetivo del equipo es solo aprobar la misión con el menor riesgo (escenario
\emph{seguridad primero}); (2) si los ensayos muestran $P_{F,c}$ por encima de
\SIrange{3.5}{7}{\percent}; (3) si las caídas muestran velocidades medias de la etapa~2
por encima de unos \SI{7.5}{m/s}, que dejan sin margen a \Req{C4}; y (4) si el calendario
no deja tiempo para la campaña de ensayos del rotor. En cualquiera de esos casos, el
análisis respalda cambiar sin dudar.

\paragraph{Por qué el equipo elige el autogiro.} El \Req{C4} no exige un autogiro: pide
``otro sistema de descenso''. El autogiro se elige por tres razones que el análisis
hace explícitas. Primera: el objetivo declarado del proyecto es desarrollar un rotor libre
y aprender de él, en línea con las misiones de autogiro de la competencia
internacional~\cite{cansat2019,cansat2025}; eso justifica $W_7\approx0{,}24$--$0{,}30$.
Segunda: con esos pesos, los dos métodos lo eligen o lo empatan, y TOPSIS lo prefiere
en el \SI{90}{\percent} de los pesos muestreados. Tercera: el autogiro cumple
\Req{C4} con probabilidad alta si funciona ($P_{\mathrm{in}}\approx1{,}00$; el riesgo
está en que funcione, no en la velocidad) y deja el eje libre para la cuerda del drogue.

\paragraph{Cómo se gestiona el riesgo de la elección.} La decisión es reversible a bajo
costo si se mantiene el paracaídas de Ø\,\SI{0.67}{m} como \textbf{plan~B}: pesa unos
\SI{50}{g}, que entran en el lastre de la configuración~A (\SI{83}{g}), y usa la misma
traba con servo. El punto de decisión es la campaña de caídas: el autogiro sigue si
alcanza al menos 15 arranques seguidos sin falla y una velocidad media de la etapa~2 de
\SI{7.5}{m/s} o menos. Con 15 éxitos en 15 intentos, el límite superior de
$P_{F}$ al \SI{60}{\percent} de confianza es $1-0{,}4^{1/15}\approx\SI{6}{\percent}$, justo
bajo el umbral de la figura~\ref{fig:decision-umbral} con los pesos de la misión técnica.
Es una evidencia débil: al \SI{90}{\percent} de confianza el mismo resultado solo asegura
$P_F<1-0{,}1^{1/15}\approx\SI{14}{\percent}$, y para demostrar $P_F<\SI{6.8}{\percent}$
con esa confianza harían falta 33 arranques seguidos sin falla. Las caídas de prueba no
son, además, idénticas al vuelo real.

\begin{nota}[Límites de este análisis]
\begin{itemize}
  \item Las $P_F$ son estimaciones; el resultado del AHP base cambia con solo pasar la
        del autogiro de \SI{4}{\percent} a \SI{3.5}{\percent}.
  \item El peso del valor técnico ($0{,}243$) es mayor que el de la tabla~\ref{tab:pugh}
        ($0{,}15$), y es el juicio que más influye en el resultado.
  \item El modelo del parapente usa coeficientes genéricos de parapentes de tamaño
        real; a esta escala ($Re\sim10^5$) su $C_L$ puede ser bastante menor.
  \item La velocidad del monocóptero y la del coaxial usan $\CR$ supuestos, sin
        ensayos propios.
  \item El movimiento angular del autogiro (\SI{25}{\degree/s}) es una conjetura; puede ser
        mayor si la fricción del cubo es alta. Se mide en las caídas con la IMU.
  \item Los juicios de K6 y K7 son del equipo. Otro equipo con otros objetivos puede
        llegar, con este mismo análisis, a elegir el paracaídas, y estaría bien.
\end{itemize}
\end{nota}

Los scripts \texttt{analysis/decision\_alternativas.py} y
\texttt{analysis/decision\_ahp.py} reproducen todas las cifras del capítulo y escriben
los datos de las figuras en \texttt{analysis/data/decision\_*.dat}.
